\subsection{Development evidence for checkpoint assessment}

Cohort A compares two development checkpoints using the same eight validation episodes. The 250-step model has 50 warmup and 200 adaptation steps; the 1,000-step model has 50 warmup and 950 adaptation steps. Mean candidate regret falls from $5.6138\times10^{-7}$ to $1.4751\times10^{-7}$, while mean closed-loop cost rises from 0.0299374 to 0.0480506. Five of six nominal episodes worsen, one improves, and two initially hot episodes tie. Both zero-delay comparisons have 4,800 requests and no misses. The learning-rate schedule also changes, so the result concerns these checkpoints rather than an isolated effect of training duration.

The absolute regret scale arises from subtracting costs for neighboring actions held for two seconds. It is neither a bound on a 60-second feedback cost nor evidence that solver uncertainty is negligible. No cost-gap-based tolerance or convergence estimate was established for the teacher candidate costs. The actual-plant refinement in Section~\ref{sec:numerical_results} leaves that teacher calculation unchanged. The differing offline and deployed admissible sets further limit an explanation based solely on future state visitation.

Cohort B contains three candidates for each of five scorer families. Two rollout-adaptation rounds each add 100 steps and 192 rows from eight training episodes, with 24 anchors per episode. Batches mix original and accumulated rollout data equally. Each policy visits its own states; shared families reach 450 total steps and the small model reaches 500. All 15 candidates pass the offline prefilter: selection score no worse than the untrained baseline, zero FV-infeasible choices, and zero unexpected fallback. The rollout rule then prioritizes adverse-event comparisons against PI and PINO, worst cost ratio over two references and five delay conditions, mean ratio, and earlier round.

Minimum-regret and rollout-based criteria select different checkpoints in every family. Normal validation cost reductions relative to the regret-minimizing choice range from 5.09\% to 68.17\%, while regret increases in every selected choice. These reductions reuse the selection data: six nominal episodes under five delays yield 30 trajectories, not 30 distinct episodes. A sensor-only regret-minimizing checkpoint produces new limit violations in 11 of 40 validation trajectories spanning three episodes, despite zero offline feasibility failures.

Only selected policies receive the separate held-out evaluation. Selected Shared-Physics has delayed normal cost 0.0146588 versus 0.0138169 for PINO-MPC and 0.0133467 for PI, with no deadline misses. The archive lacks an independent comparison of regret-selected and rollout-selected alternatives. B therefore motivates checking the selection criterion but does not validate a selector's test benefit. It is also distinct from D, whose two residual systems share training examples and batch order.

